\chapter{What Constructive Logic Refuses to Guess}
\label{ch:constructive-boundary}

The constructive core has a simple discipline: to prove a disjunction, construct a branch; to prove an existential, supply a witness. Classical logic validates principles that need not expose such information. The difference is not a defect in either system; it is a difference in admissible proof principles.

Let $\neg A$ abbreviate $A\to\Empty$. The schema $A+\neg A$ is excluded middle in the propositional reading. It is not derivable from the pure $\mathrm{CH}\text{-}0$ rules alone.

\begin{figure}[p]\centering\begin{tikzpicture}[gclplate]\node[anchor=west,font=\sffamily\bfseries\Large,gcllabel] at (-6,3.8){A WORLD WHERE EXCLUDED MIDDLE FAILS};\draw[GCLWarm,line width=1pt](-6,3.45)--(6,3.45);\node[gcllabel,draw](w0) at (-3,1.2){$w_0$: not $P$};\node[gcllabel,draw](w1) at (3,1.2){$w_1$: forces $P$};\draw[gclwarmarrow](w0)--node[above,gcllabel]{$w_0\le w_1$}(w1);\node[gcllabel] at (0,-.2){$w_0\not\Vdash P\lor\neg P$};\end{tikzpicture}\caption{\textbf{Plate 22. Finite Kripke countermodel.} At the root, neither $P$ nor $\neg P$ is forced, so excluded middle fails there. Scope: selected intuitionistic Kripke semantics.}\label{plate:kripke-em}\end{figure}


To make that nonderivability visible, consider the two-world Kripke frame $w_0\le w_1$. Let atomic $P$ be false at $w_0$ and forced at $w_1$. Persistence holds because once $P$ becomes true it remains true upward.

At $w_0$, $P$ is not forced. Neither is $\neg P$: forcing $P\to\Empty$ at $w_0$ would require every extension forcing $P$ also to force $\Empty$, but $w_1$ forces $P$ and not $\Empty$. Therefore $w_0$ forces neither disjunct, so it does not force $P+\neg P$.

\begin{proposition}[Finite countermodel to excluded middle]
The Kripke model above does not force $P\lor\neg P$ at its root.
\end{proposition}
\begin{proof}
By the forcing clauses for disjunction and implication just evaluated: neither $P$ nor $P\to\Empty$ is forced at $w_0$.
\end{proof}

By the standard soundness theorem for intuitionistic propositional logic with respect to Kripke semantics, the existence of this countermodel implies that excluded middle is not derivable in the corresponding intuitionistic calculus.

\begin{figure}[p]\centering\begin{tikzpicture}[gclplate]\node[anchor=west,font=\sffamily\bfseries\Large,gcllabel] at (-6,3.8){CONSTRUCTIVE INFORMATION};\draw[GCLWarm,line width=1pt](-6,3.45)--(6,3.45);\node[gcllabel,draw] at (-3,1){construct branch + witness};\node[gcllabel,draw] at (3,1){classical validity may omit direct branch data};\end{tikzpicture}\caption{\textbf{Plate 23. Constructive information versus classical commitment.} Constructive introduction rules expose branch/witness data that classical validity need not provide in the same form. Scope: calculus-relative, not a value judgment.}\label{plate:constructive-classical}\end{figure}


\begin{workedexample}[title={Information not supplied}]
A classical proof may establish that one of $P$ or $\neg P$ must hold without constructing which branch in the same direct manner as $+I_1$ or $+I_2$. If we want executable content from such reasoning, we need a translation, additional semantic structure, or control behavior. Chapter~\ref{ch:classical-control} explores one such route.
\end{workedexample}

\begin{figure}[p]\centering\begin{tikzpicture}[gclplate]\node[anchor=west,font=\sffamily\bfseries\Large,gcllabel] at (-6,3.8){DOUBLE NEGATION IS A BOUNDARY};\draw[GCLWarm,line width=1pt](-6,3.45)--(6,3.45);\node[gcllabel,draw](c) at (-3,1.2){classical source};\node[gcllabel,draw](n) at (3,1.2){negative/CPS target};\draw[gclbluearrow](c)--node[above,gcllabel]{translation}(n);\end{tikzpicture}\caption{\textbf{Plate 24. Double negation as a translation boundary.} Classical reasoning can enter a constructive target through a negative translation rather than literal proof identity. Scope: selected translation only.}\label{plate:double-negation-boundary}\end{figure}


\begin{warningbox}
“Nonconstructive” is calculus-relative. It does not mean invalid, useless, or irrational. It says that the proof uses principles not justified by the witness-producing rules fixed for $\mathrm{CH}\text{-}0$.
\end{warningbox}

\begin{labbox}
Run \texttt{python labs/lab08_kripke.py}. The finite model checker verifies persistence and evaluates the selected excluded-middle formula at both worlds. Its result is a countermodel for that formula under the implemented forcing semantics, not a survey of all constructive logics.
\end{labbox}

\section*{Exercises}\addcontentsline{toc}{section}{Exercises}
\begin{enumerate}[label=\textbf{8.\arabic*.}]
\item \textbf{Checkpoint.} Define $\neg A$ in $\mathrm{CH}\text{-}0$.
\item \textbf{Checkpoint.} What evidence does constructive disjunction require?
\item \textbf{Checkpoint.} Why does $w_0$ fail to force $\neg P$ in the example?
\item \textbf{Core.} Evaluate $P\lor\neg P$ at $w_1$.
\item \textbf{Core.} State the persistence condition for atomic forcing.
\item \textbf{Core.} Build a finite model separating another selected classical schema from intuitionistic validity.
\item \textbf{Synthesis.} Explain how a semantic countermodel becomes a syntactic nonderivability result only with a soundness theorem.
\item \textbf{Synthesis.} Contrast lack of a constructive proof with falsity in a classical two-valued model.
\item \textbf{Proof Workshop.} Prove persistence for the implication forcing clause by induction on formulas.
\item \textbf{Proof Workshop.} Verify every forcing clause used in the root countermodel.
\item \textbf{Design Clinic.} Specify checks that a finite Kripke-model parser should impose before evaluation.
\item \textbf{Challenge.} Explain why adding excluded middle changes proof theory even when ordinary classical truth tables remain familiar.
\end{enumerate}
